%AttackName: M-ZOO TO-7

%Input:
\textbf{$\mathbf{x}$}: Original input image. \\
\textbf{$\delta$}: Perturbation to be found. \\
\textbf{$C(\mathbf{x} + \delta)$}: Constraint requiring the modified image to be classified as a specific class. \\
\textbf{$\epsilon$}: Perturbation bound. \\
\textbf{$\text{model}$}: The target model for the attack.

%Output:
The optimized perturbation $\delta_{\text{opt}}$ that minimizes the loss function while adhering to the perturbation constraint.

%Formula:
The ZOO attack Targeted One-shot Optimization optimizes the perturbation $\delta$ as follows:
\begin{equation*}
\begin{aligned}
\min_{\delta} & \quad \mathcal{L}(\mathbf{x} + \delta) \\
\text{s.t.} & \quad \|\delta\|_p \leq \epsilon
\end{aligned}
\end{equation*}

where $\mathcal{L}$ is the loss function based on the model's output.

The perturbation $\delta$ can be estimated using the zeroth-order optimization:
\begin{equation*}
\delta_i \approx \frac{f(\mathbf{x} + \epsilon \mathbf{e}_i) - f(\mathbf{x} - \epsilon \mathbf{e}_i)}{2 \epsilon} \mathbf{e}_i
\end{equation*}

where $f$ is the output of the classifier, and $\mathbf{e}_i$ is the unit vector in the direction of the $i$-th dimension.

%Explanation:
The ZOO attack uses a zeroth-order optimization method to estimate gradients, allowing it to generate adversarial examples without requiring access to the model's gradients. This makes it effective against black-box models while ensuring the perturbation stays within the specified bound.